\section{Activity Method}

A second survey recorded when ground beetles were most active.

Beetles were sampled along hedgerows.

Species were identified from published keys.

Only one season was sampled, and the pattern may differ in other years.

\section{Activity Results}

Beetle activity peaked in June, with 41 percent of all captures. Activity fell below 5 percent per month after August. Beetles were therefore most active in early summer.

We were delighted to find Carabus auratus in the third hedgerow.

The timing of activity showed no link to hedgerow width (F = 0.8, p = .45). Width thus appeared to shape how many species a plot held and to leave their season of activity unchanged.

\section{Discussion}

Wider hedgerows hold more species, which suggests that width adds habitat. The result extends earlier work on hedgerows as shelter for beetles (Lang, 2016) to the width of the hedgerow.

The pattern matches earlier surveys in Denmark (Holm, 2012).

Older hedgerows are wider in this sample, so age may explain part of the gain. A survey of hedgerow age on the same plots would separate the two effects.

A 5 m strip removes land from the crop; the same strip holds 16 species on average. Farmers therefore weigh a loss of cropland against a gain in species.

The gain relates to width, not to hedgerow length, since every transect was 20 m long.

We suggest that agri-environment schemes consider payments for hedgerows wider than 5 m.

\section{Conclusion}

Hedgerows wider than 5 m may hold more ground beetle species on arable farms. Future work should follow the same plots across at least three seasons.
